% =============================================================
% Lancet-style two-column scientific paper
% Variabilidad Hospitalaria Oncológica — UDD 2026
% =============================================================
\documentclass[10pt,twocolumn]{article}

% ── Encoding & language ──────────────────────────────────────
\usepackage[utf8]{inputenc}
\usepackage[T1]{fontenc}
\usepackage[english]{babel}

% ── Fonts (close to Lancet's Minion/Myriad style) ────────────
\usepackage{mathptmx}          % Times-like serif for body text
\usepackage[scaled=0.90]{helvet}
\renewcommand{\familydefault}{\rmdefault}

% ── Page geometry ────────────────────────────────────────────
\usepackage[
  top=2.0cm, bottom=2.2cm,
  left=1.6cm, right=1.6cm,
  columnsep=0.6cm
]{geometry}

% ── Core packages ────────────────────────────────────────────
\usepackage{graphicx}
\usepackage{booktabs}
\usepackage{array}
\usepackage{tabularx}
\usepackage{longtable}
\usepackage{multirow}
\usepackage{xcolor}
\usepackage{mdframed}
\usepackage{setspace}
\usepackage{parskip}
\usepackage{enumitem}
\usepackage{microtype}
\usepackage{caption}
\usepackage{amsmath}
\usepackage{amssymb}
\usepackage{url}
\usepackage{hyperref}
\usepackage{fancyhdr}
\usepackage{titlesec}
\usepackage{abstract}
\usepackage{authblk}
\usepackage{balance}
\usepackage{flushend}
\usepackage{lipsum}

% ── Lancet colour palette ─────────────────────────────────────
\definecolor{lancetgreen}{RGB}{0,96,70}      % Lancet dark green
\definecolor{lancetlightgreen}{RGB}{0,130,95}
\definecolor{lancetgray}{RGB}{245,245,245}
\definecolor{lancetborder}{RGB}{0,96,70}
\definecolor{boxbg}{RGB}{237,246,242}
\definecolor{tablehead}{RGB}{0,96,70}

% ── Hyperref setup ────────────────────────────────────────────
\hypersetup{
  colorlinks=true,
  linkcolor=lancetgreen,
  citecolor=lancetgreen,
  urlcolor=lancetgreen,
  pdftitle={Hospital-Level Variation in Oncological Treatment...},
  pdfauthor={Rodriguez, Amat, Herrera}
}

% ── Section title formatting ──────────────────────────────────
\titleformat{\section}
  {\normalfont\bfseries\color{lancetgreen}\fontsize{10}{12}\selectfont}
  {}{0pt}{}[\vspace{-2pt}\color{lancetgreen}\titlerule[0.6pt]]

\titleformat{\subsection}
  {\normalfont\bfseries\fontsize{9.5}{11}\selectfont\itshape}
  {}{0pt}{}

\titlespacing*{\section}{0pt}{8pt}{4pt}
\titlespacing*{\subsection}{0pt}{5pt}{2pt}

% ── Caption styling ──────────────────────────────────────────
\captionsetup{
  font={small},
  labelfont={bf,color=lancetgreen},
  labelsep=period,
  justification=justified,
  singlelinecheck=false
}

% ── Table styling ─────────────────────────────────────────────
\newcolumntype{L}[1]{>{\raggedright\arraybackslash}p{#1}}
\newcolumntype{C}[1]{>{\centering\arraybackslash}p{#1}}
\newcolumntype{R}[1]{>{\raggedleft\arraybackslash}p{#1}}

% ── Research-in-context box ──────────────────────────────────
\mdfdefinestyle{ricbox}{
  backgroundcolor=boxbg,
  linecolor=lancetgreen,
  linewidth=1.5pt,
  topline=true, bottomline=true, leftline=true, rightline=true,
  innertopmargin=8pt, innerbottommargin=8pt,
  innerleftmargin=10pt, innerrightmargin=10pt,
  skipabove=10pt, skipbelow=10pt
}

% ── Header/footer ────────────────────────────────────────────
\pagestyle{fancy}
\fancyhf{}
\renewcommand{\headrulewidth}{0pt}
\fancyhead[L]{\footnotesize\textit{Articles}}
\fancyhead[R]{\footnotesize\textcolor{lancetgreen}{\textit{www.thelancet.com}}}
\fancyfoot[C]{\footnotesize\thepage}

% ── Abstract environment ─────────────────────────────────────
\renewenvironment{abstract}
  {\small\setlength{\parindent}{0pt}\setlength{\parskip}{3pt}%
   \noindent\rule{\linewidth}{1.5pt}\vspace{4pt}\par}
  {\par\vspace{4pt}\noindent\rule{\linewidth}{0.8pt}}

% ── Footnote size ────────────────────────────────────────────
\renewcommand{\footnotesize}{\fontsize{8}{9.5}\selectfont}
\setlength{\footnotesep}{4pt}

% ── Line spacing ─────────────────────────────────────────────
\setstretch{1.08}

% ── Reference list style ─────────────────────────────────────
\setlist[enumerate,1]{label=\arabic*.,leftmargin=16pt,
  itemsep=0pt,parsep=0pt,topsep=2pt,font=\small}

% =============================================================
\begin{document}

% ── TITLE BLOCK (full width) ─────────────────────────────────
\twocolumn[{%
\begin{@twocolumnfalse}

\noindent\colorbox{lancetgreen}{%
  \parbox{\dimexpr\linewidth}{\vspace{4pt}
    \hfill\textcolor{white}{\fontsize{7.5}{9}\selectfont
      \textbf{Articles} \quad|\quad
      \textit{Análisis de Datos e Inferencia Estadística — Universidad del Desarrollo, 2026}}
    \hspace{6pt}\vspace{2pt}
  }
}

\vspace{10pt}

{\fontsize{15}{19}\selectfont\bfseries
Hospital-Level Variation in Gastric Cancer Treatment
and Its Effects on Mortality and Length of Stay
in the Chilean Public Health System\par}

\vspace{8pt}

{\fontsize{9}{11}\selectfont
Vicente Rodríguez,$^{a,*}$ José Tomás Amat,$^{a}$ Sebastián Herrera$^{a}$\par}

\vspace{6pt}
{\fontsize{8}{10}\selectfont\itshape
$^{a}$School of Engineering, Universidad del Desarrollo, Santiago, Chile\\[2pt]
$^{*}$Correspondence: vjrodrig2004@gmail.com\par}

\vspace{8pt}

% ── ABSTRACT ─────────────────────────────────────────────────
\begin{abstract}
\noindent\textbf{\textcolor{lancetgreen}{Background}}\quad
Gastric cancer is the leading cause of cancer-related mortality in the Chilean public health system,
yet the extent to which hospital characteristics—rather than patient clinical profiles—determine
outcomes remains unquantified. We investigated whether the treating hospital explains
significant variation in length of stay, number of procedures, and in-hospital mortality
in gastric cancer patients, after controlling for clinical severity.

\vspace{4pt}
\noindent\textbf{\textcolor{lancetgreen}{Methods}}\quad
We performed a cross-sectional analysis of the national Diagnosis-Related Groups (DRG)
database covering 100\% of hospitalizations in public FONASA-covered institutions between
2019 and 2024. We identified 17\,335 gastric cancer discharges (ICD-10: C16.\,*) and
restricted regression analyses to 9\,855 discharges across 15 high-volume hospitals.
We applied Kruskal-Wallis tests with Dunn-Bonferroni post-hoc correction (H\textsubscript{1}),
multivariable logistic regression (H\textsubscript{2}), and multiple OLS with HC3 robust
standard errors (H\textsubscript{3}). A random-intercept mixed-effects model (MixedLM) was
used to estimate the intraclass correlation coefficient (ICC). Socioeconomic context
was incorporated through linkage with 2024 Municipal Human Development Index (HDI)
quintiles. Hospital typology was characterized using K-means clustering.

\vspace{4pt}
\noindent\textbf{\textcolor{lancetgreen}{Findings}}\quad
The distribution of procedures differed significantly across hospitals (Kruskal-Wallis
$H = 412.3$, $p < 0{.}001$, $\varepsilon^2 = 0{.}08$), with over 60\% of hospital pairs
showing significant pairwise differences after Bonferroni correction. Each additional
procedure increased the odds of in-hospital mortality by 5\%
(OR\,=\,1{.}05; 95\%\,CI 1.02–1.08; $p < 0{.}001$), while GRD severity was the
dominant predictor (OR\,=\,6{.}1). Each additional procedure was also associated with
a 9{.}7\% increase in length of stay ($\beta = 0{.}093$; $R^2 = 0{.}636$). The
ICC from the multilevel model indicated that 18{.}4\% of the variance in log-transformed
length of stay is attributable to the hospital rather than to the patient. Patients
residing in the lowest HDI quintile (Q1) showed 2{.}8 percentage points higher
in-hospital mortality than those in the highest quintile (Q5) ($p = 0{.}003$).
K-means clustering identified three distinct hospital typologies with significant
differences in resource utilization and patient outcomes.

\vspace{4pt}
\noindent\textbf{\textcolor{lancetgreen}{Interpretation}}\quad
Hospital-level variation in gastric cancer care in Chile is structural, statistically
robust, and socioeconomically patterned. The establishment of care explains nearly one fifth
of the variance in length of stay after controlling for patient characteristics. These
findings support policies targeting institutional standardization of oncological protocols
and differential resource allocation according to territorial socioeconomic disadvantage.

\vspace{6pt}
{\fontsize{8}{10}\selectfont
\textit{Keywords:} gastric cancer, hospital variation, intraclass correlation, Chile,
DRG database, health equity, socioeconomic determinants\par}
\end{abstract}

\vspace{10pt}
\end{@twocolumnfalse}
}]

% ── RESEARCH IN CONTEXT BOX ──────────────────────────────────
\begin{mdframed}[style=ricbox]
{\fontsize{8.5}{11}\selectfont\bfseries Research in context\par}
\vspace{4pt}
{\fontsize{8}{10}\selectfont
\textbf{Evidence before this study}\\
Hospital-level variation in oncological outcomes has been documented in high-income
countries (Kamaraju et al., 2022; OECD, 2023). However, evidence from Latin America
using national administrative data is scarce. Prior work on Chile's GRD database has
focused on multimorbidity prevalence (Bernal et al., 2026) but not on inter-institutional
variability in treatment intensity or its socioeconomic determinants.

\vspace{4pt}
\textbf{Added value of this study}\\
This is the first nationwide analysis of hospital-level variation in gastric cancer care in
Chile using DRG data covering six years (2019–2024) and 15 high-volume public hospitals.
We quantify the institutional contribution to variance in outcomes through a random-intercept
mixed model ICC (18.4\%), link treatment variation to territorial socioeconomic deprivation
via HDI quintiles, and characterize three distinct hospital typologies through unsupervised
clustering—providing actionable evidence for health system governance.

\vspace{4pt}
\textbf{Implications of all available evidence}\\
The hospital of care matters beyond what can be explained by the patient's clinical profile.
This institutional effect is unevenly distributed: patients from lower-HDI municipalities
disproportionately access hospitals in the high-complexity/urgent cluster, closing a
cycle of disadvantage. Standardization of oncological protocols and equitable resource
distribution are the priority interventions identified.
}
\end{mdframed}

% ═══════════════════════════════════════════════════
\section{Introduction}
% ═══════════════════════════════════════════════════

Cancer is the second leading cause of death in Chile and the primary contributor to
years of life lost prematurely.\textsuperscript{1} Among digestive cancers, gastric
cancer (ICD-10: C16) ranks first in incidence and mortality within the public health
system, with five-year survival rates substantially lower than those reported in
high-income countries.\textsuperscript{2}

A growing body of international evidence demonstrates that hospital-level variation in
cancer care—differences in procedure use, treatment intensity, and outcomes across
institutions—is not fully explained by the clinical characteristics of admitted
patients.\textsuperscript{3,4} This \textit{unjustified clinical variability}
(Wennberg, 2010)\textsuperscript{5} has direct implications for health equity: when
outcomes are partly determined by where a patient receives care rather than by their
clinical needs, access to quality treatment becomes a function of geography and
socioeconomic status rather than medical necessity.

In Chile, the public health system (FONASA) covers approximately 80\% of the
population and operates through a network of hospitals with heterogeneous levels of
complexity, resource endowment, and geographic distribution. Hospitals serving
lower-income territories—identified through the Municipal Human Development Index
(HDI)—may face additional barriers, including higher proportions of late-stage
presentations, emergency admissions, and constrained treatment capacity. The
intersection of institutional variation and socioeconomic disadvantage thus creates
a compound risk for patients in the most vulnerable communities.

Despite the availability of a nationwide DRG database covering all public
hospitalizations, no study has quantified inter-hospital variation in gastric cancer
outcomes using multilevel methods, nor linked this variation to the socioeconomic
profile of the territory of care. We address this gap by testing three pre-specified
hypotheses: (H\textsubscript{1}) the distribution of procedures differs significantly
across hospitals; (H\textsubscript{2}) procedure count is associated with in-hospital
mortality after adjusting for clinical covariates; and (H\textsubscript{3}) procedure
count predicts length of stay after similar adjustment. We further quantify the
institutional contribution to variance through the ICC, characterize the
socioeconomic gradient in outcomes, and identify distinct hospital typologies through
unsupervised clustering.

% ═══════════════════════════════════════════════════
\section{Methods}
% ═══════════════════════════════════════════════════

\subsection{Data source and study population}

Data were obtained from Chile's national Diagnosis-Related Groups (DRG) database,
maintained by the Ministry of Health (MINSAL) and administered through FONASA.
This database compiles 100\% of hospital discharge records from all public institutions
providing care to FONASA-insured patients—approximately 80\% of the Chilean
population—across all regions and levels of hospital complexity. We analyzed discharges
recorded between January 2019 and December 2024 (six calendar years).

We selected all discharges with a primary diagnosis of malignant gastric neoplasm
(ICD-10: C16 and subcategories, ``gastric cancer'') as the focal population. Five
sequential cleaning stages were applied: (1) restriction to oncological ICD-10 codes
C00–D49; (2) exclusion of obstetric admissions (TIPO\_INGRESO\,=\,OBSTÉTRICA),
which represent a clinically incomparable process; (3) truncation of length-of-stay
outliers at the 99th percentile ($P_{99} = 25$\,days); (4) restriction to C16.*
subcodes; and (5) further restriction to the 15 highest-volume hospitals (each with
$\geq 30$ C16.* discharges) for regression analyses to ensure model convergence
(Fig.~S1, appendix).

\subsection{Variables}

\textit{Primary outcomes.} (i) \textbf{Quantity of procedures} (\texttt{cantidad\_pro\-ce\-di\-mien\-tos}): count of non-null procedural fields per discharge record;
(ii) \textbf{In-hospital mortality} (\texttt{mortalidad\_int}): binary indicator
(1\,=\,discharge type\,=\,deceased); (iii) \textbf{Length of hospital stay}
(\texttt{dias\_estada}): days from admission to discharge, log-transformed
[$\texttt{log\_dias} = \log(1 + \texttt{dias\_estada})$] to correct right skew
(skewness: 1.63 raw vs.\ $-$0.32 after transformation).

\textit{Clinical covariates.} Age at hospitalization; GRD severity score
(\texttt{se\-ve\-ri\-dad\_grd}, 1–4 ordinal scale representing clinical complexity);
GRD relative weight (\texttt{pe\-so\_grd}, continuous proxy for resource consumption);
comorbidity count (number of active secondary diagnoses across DIAG2–DIAG35 fields,
excluding the primary diagnosis—a simplified Charlson-Deyo proxy).

\textit{Institutional variable.} Hospital identifier, entered as a fixed effect
(regression models) or random effect (multilevel model).

\textit{Socioeconomic variable.} Municipal Human Development Index (HDI), 2024 edition,
sourced externally and linked to the patient's municipality of residence through
normalized name-matching. HDI was categorized into five quintiles (Q1 = lowest to
Q5 = highest) based on the distribution observed in the analytic sample.

\subsection{Statistical analysis}

\textbf{Exploratory data analysis.} Continuous variables were summarized as mean\,$\pm$\,SD
and median (IQR); categorical variables as frequencies and percentages. Pearson
correlation matrices and bivariate cross-tabulations were constructed for all
outcome–covariate pairs.

\textbf{H\textsubscript{1} — Kruskal-Wallis + Dunn post-hoc.}
Normality of procedure counts was assessed via Shapiro-Wilk test on a stratified
subsample ($n = 5{,}000$). Given rejection of normality ($W = 0{.}922$,
$p < 0{.}001$), a Kruskal-Wallis rank-sum test was applied across the top-25
hospitals by volume ($n \geq 30$ discharges each). Significant omnibus results were
followed by pairwise Dunn-Bonferroni post-hoc tests. Effect size was quantified
with epsilon-squared ($\varepsilon^2$).

\textbf{H\textsubscript{2} — Logistic regression (in-hospital mortality).}
A multivariable logistic regression was specified as:
\begin{equation*}
  \text{logit}[\Pr(\text{mortality})] = \beta_0 + \beta_1\,\text{proc} +
  \beta_2\,\text{age} + \beta_3\,\text{sev} + \beta_4\,\text{wt} +
  \beta_5\,\text{comorb} + \boldsymbol{\gamma}\,C(\text{hospital})
\end{equation*}
estimated via BFGS with 500 iterations (statsmodels v0.14). Discrimination was
assessed by AUC-ROC on a stratified 75/25 train-test split. Odds ratios (OR) are
presented with 95\% confidence intervals (CI). Class imbalance (mortality rate\,$\approx$\,4\%)
was addressed through three strategies: baseline, class-weight balancing, and
SMOTE oversampling, compared on sensitivity at threshold 0.10.

\textbf{H\textsubscript{3} — Multiple OLS regression (length of stay).}
An OLS model was specified as:
\begin{equation*}
  \log(1 + \text{LOS}) = \beta_0 + \beta_1\,\text{proc} + \beta_2\,\text{age} +
  \beta_3\,\text{sev} + \beta_4\,\text{wt} + \beta_5\,\text{comorb} +
  \boldsymbol{\gamma}\,C(\text{hospital}) + \varepsilon
\end{equation*}
estimated with HC3 heteroskedasticity-robust standard errors (White, 1980;
MacKinnon \& White, 1985). Back-transformation to the original scale used
$\hat{y}_{\text{orig}} = \exp(\hat{y}_{\log}) - 1$ (\texttt{numpy.expm1}).
Predictive performance was evaluated on the hold-out test set (MAE, RMSE, $R^2$).

\textbf{Multilevel model (ICC).} A random-intercept linear mixed model was fitted
with \texttt{statsmodels MixedLM} using restricted maximum likelihood (REML).
The ICC was computed as:
\begin{equation*}
  \text{ICC} = \frac{\sigma^2_{\text{hospital}}}{\sigma^2_{\text{hospital}} +
  \sigma^2_{\text{residual}}}
\end{equation*}
Hospital-level random effects (BLUPs) were extracted and visualized in a forest plot.

\textbf{Socioeconomic analysis.} Kruskal-Wallis tests compared length of stay and
chi-squared tests compared in-hospital mortality across HDI quintiles. Standardized
residuals from the contingency tables identified which quintiles drove significant
deviations.

\textbf{Hospital typology (K-means clustering).} Hospital profiles were constructed
aggregating five indicators: in-hospital mortality rate, mean length of stay,
mean procedure count, emergency admission rate, and mean GRD severity. Features
were standardized (z-scores) and cluster number was selected via the elbow method
and silhouette coefficient ($k = 3$). Principal Component Analysis (two components)
was used for two-dimensional visualization.

All analyses were conducted in Python 3.13 (pandas 2.x, numpy, statsmodels 0.14,
scikit-learn 1.4, matplotlib, seaborn). The significance threshold was
$\alpha = 0{.}05$ throughout. All procedures were non-experimental; no ethics
approval was required for analysis of deidentified administrative data.

% ═══════════════════════════════════════════════════
\section{Results}
% ═══════════════════════════════════════════════════

\subsection{Study population characteristics}

Between 2019 and 2024, the DRG database contained 457\,717 total hospital
discharges from FONASA-covered institutions. After applying the five-stage cleaning
pipeline, 421\,824 oncological discharges (C00–D49) were retained, of which
17\,335 corresponded to gastric cancer (C16.*), representing 4.1\% of the
oncological universe. The regression analytic sample comprised 9\,855 discharges
across 15 public hospitals after exclusion of low-volume establishments
(Table~\ref{tab:descriptive}).

The mean patient age was $63.2 \pm 13.1$ years. Male patients accounted for
60.3\% of C16.* discharges. The median length of stay was 7\,days (IQR 3–12),
with right skew (skewness\,=\,1.63) justifying logarithmic transformation. The
median number of procedures per admission was 3 (IQR 1–6). The overall
in-hospital mortality rate was 4{.}05\%, consistent with a clinically severe,
late-stage presentation profile.

\begin{table}[t]
\small
\begin{tabularx}{\linewidth}{@{}L{4.5cm}R{1.5cm}R{1.5cm}@{}}
\toprule
\rowcolor{tablehead}
\textcolor{white}{\textbf{Characteristic}} &
\textcolor{white}{\textbf{C00–D49}} &
\textcolor{white}{\textbf{C16.*}} \\
\midrule
\textit{n} (discharges)                     & 421\,824    & 9\,855$^*$  \\
Age, mean\,$\pm$\,SD (years)                & 59.8\,$\pm$\,15.4 & 63.2\,$\pm$\,13.1 \\
Male sex, \%                                & 48.1        & 60.3       \\
Length of stay, median [IQR] (days)         & 5 [2–10]    & 7 [3–12]   \\
Procedures, median [IQR]                    & 2 [1–4]     & 3 [1–6]    \\
GRD severity, mean\,$\pm$\,SD              & 1.82\,$\pm$\,0.91 & 2.14\,$\pm$\,0.88 \\
GRD weight, mean\,$\pm$\,SD                & 1.31\,$\pm$\,1.02 & 1.68\,$\pm$\,1.24 \\
Comorbidities, mean\,$\pm$\,SD             & 1.43\,$\pm$\,1.78 & 1.61\,$\pm$\,1.84 \\
Emergency admission, \%                     & 38.9        & 42.7       \\
In-hospital mortality, \%                   & 2.31        & 4{.}05     \\
\bottomrule
\end{tabularx}
\caption{Descriptive statistics of the study population. $^*$Restricted to top-15
hospitals ($\geq$30 C16.* discharges each). C00–D49: all oncological discharges;
C16.*: gastric cancer. IQR = interquartile range; SD = standard deviation.}
\label{tab:descriptive}
\end{table}

\subsection{H\textsubscript{1}: Kruskal-Wallis — procedural variability}

The Shapiro-Wilk test rejected normality of procedure counts
($W = 0{.}922$, $p < 0{.}001$), justifying non-parametric testing. The Kruskal-Wallis
test across the top-25 hospitals ($N = 14{,}283$; $k = 25$ hospitals) yielded
$H(24) = 412{.}3$, $p < 0{.}001$, with $\varepsilon^2 = 0{.}080$, corresponding to
a medium-to-large effect. Dunn-Bonferroni post-hoc comparisons identified
significant pairwise differences in 61{.}8\% of all hospital pairs (Fig.~1).

This pattern indicates that procedural variability in gastric cancer care is
\textit{structural and generalized}: it is not driven by a single outlier institution
but reflects system-wide heterogeneity in treatment intensity that persists
after controlling for case mix through the GRD classification.

\subsection{H\textsubscript{2}: Logistic regression — in-hospital mortality}

The multivariable logistic model was fitted on the full sample ($n = 9{,}855$)
and evaluated on a stratified hold-out set ($n = 2{,}464$). The model achieved
an AUC-ROC of $0{.}823$ (95\% bootstrap CI: 0.80–0.84), indicating good
discriminative ability (Table~\ref{tab:or}).

\begin{table}[t]
\small
\begin{tabularx}{\linewidth}{@{}L{3.8cm}C{1.0cm}C{1.6cm}C{1.0cm}@{}}
\toprule
\rowcolor{tablehead}
\textcolor{white}{\textbf{Predictor}} &
\textcolor{white}{\textbf{OR}} &
\textcolor{white}{\textbf{95\% CI}} &
\textcolor{white}{\textbf{\textit{p}}} \\
\midrule
Procedures (per unit)  & 1{.}05 & 1.02–1.08 & $<$0{.}001 \\
Age (per year)         & 1{.}01 & 0.99–1.02 & 0{.}341    \\
GRD severity (per unit)& 6{.}12 & 5.42–6.91 & $<$0{.}001 \\
GRD weight (per unit)  & 0{.}88 & 0.82–0.95 & 0{.}001    \\
Comorbidity (per count)& 1{.}06 & 1.01–1.11 & 0{.}018    \\
Hospital (fixed effects)& \multicolumn{3}{c}{\textit{included, not shown}} \\
\midrule
Pseudo-$R^2$ (McFadden) & \multicolumn{3}{c}{0{.}182}  \\
AUC-ROC                 & \multicolumn{3}{c}{0{.}823}  \\
$N$                     & \multicolumn{3}{c}{9\,855}   \\
\bottomrule
\end{tabularx}
\caption{Odds ratios from multivariable logistic regression for in-hospital
mortality, gastric cancer (C16.*). Fixed effects for hospital are included
in the model but omitted from the table for brevity. OR = odds ratio;
CI = confidence interval.}
\label{tab:or}
\end{table}

GRD severity dominated the prediction (OR\,=\,6{.}12), reflecting the expected
role of clinical complexity in mortality. The OR for procedures
(1{.}05 per additional procedure) should be interpreted cautiously due to
\textit{indication confounding}: more severely ill patients receive more procedures,
inflating the apparent procedure–mortality association even after adjustment.
Age was not a significant independent predictor after conditioning on GRD severity.
The protective effect of GRD weight (OR\,$<$\,1) reflects that higher relative
weights are assigned to complex but schedulable (elective) procedures with
better planning and outcomes.

\subsection{H\textsubscript{3}: OLS regression — length of stay}

The OLS model for $\log(1 + \text{LOS})$ achieved $R^2 = 0{.}636$
(adjusted $R^2 = 0{.}635$), $F(19,\,9{,}835) = 978{.}1$, $p < 0{.}001$
(Table~\ref{tab:ols}). In the hold-out evaluation, back-transformed predictions
yielded MAE = 3{.}47\,days and RMSE = 5{.}61\,days.

\begin{table}[t]
\small
\begin{tabularx}{\linewidth}{@{}L{3.8cm}C{0.9cm}C{0.9cm}C{1.4cm}C{0.9cm}@{}}
\toprule
\rowcolor{tablehead}
\textcolor{white}{\textbf{Predictor}} &
\textcolor{white}{\textbf{$\hat\beta$}} &
\textcolor{white}{\textbf{SE}} &
\textcolor{white}{\textbf{95\% CI}} &
\textcolor{white}{\textbf{$\Delta$\%}} \\
\midrule
Procedures (per unit)   & 0{.}093 & 0{.}004 & 0.085–0.101 & +9.7  \\
Age (per year)          & 0{.}002 & 0{.}001 & 0.001–0.004 & +0.2  \\
GRD severity (per unit) & 0{.}241 & 0{.}011 & 0.220–0.263 & +27.3 \\
GRD weight (per unit)   & 0{.}078 & 0{.}008 & 0.062–0.094 & +8.1  \\
Comorbidity (per count) & 0{.}031 & 0{.}006 & 0.020–0.043 & +3.2  \\
Hospital (fixed effects)& \multicolumn{4}{c}{\textit{14 dummies, included}} \\
\midrule
$R^2$ & \multicolumn{4}{c}{0{.}636}   \\
Adj.\,$R^2$ & \multicolumn{4}{c}{0{.}635} \\
$F(19, 9835)$ & \multicolumn{4}{c}{978{.}1, $p < 0.001$} \\
MAE (test, orig.) & \multicolumn{4}{c}{3{.}47 days} \\
RMSE (test, orig.) & \multicolumn{4}{c}{5{.}61 days} \\
$N$ & \multicolumn{4}{c}{9\,855} \\
\bottomrule
\end{tabularx}
\caption{OLS regression coefficients for $\log(1 + \text{LOS})$, HC3
robust standard errors. $\hat\beta$ = estimated coefficient;
$\Delta\%$ = approximate percentage change in LOS per unit increase
$[(\exp(\hat\beta)-1)\times 100]$. SE = standard error; CI = confidence interval.}
\label{tab:ols}
\end{table}

Each additional procedure was associated with a 9{.}7\% increase in LOS, holding
all other covariates constant. GRD severity had the largest marginal effect
(+27{.}3\% per severity level), confirming that clinical complexity is the primary
driver of hospitalization duration. Fixed effects for hospital were collectively
significant ($F = 14{.}8$, $p < 0{.}001$), indicating that institutional
characteristics explain additional variation in LOS beyond patient-level covariates.

\subsection{Multilevel model and intraclass correlation (ICC)}

The random-intercept MixedLM model (REML) estimated the following variance
components: $\hat\sigma^2_{\text{hospital}} = 0{.}041$ (between-hospital variance)
and $\hat\sigma^2_{\text{residual}} = 0{.}183$ (within-hospital variance), yielding:

\begin{equation*}
  \text{ICC} = \frac{0{.}041}{0{.}041 + 0{.}183} = 0{.}184 \quad (18{.}4\%)
\end{equation*}

This indicates that \textbf{18{.}4\% of the total variance in log-transformed LOS is
attributable to the hospital of care}, rather than to patient-level clinical
characteristics. The forest plot of hospital-level random effects (Fig.~2) reveals
that some hospitals have adjusted LOS systematically higher or lower than the
system mean for clinically equivalent patients.

Fixed-effect coefficients from the multilevel model were consistent with those
from the OLS with fixed hospital effects (within 5\% for all continuous predictors),
providing evidence that the OLS estimates are robust to the choice of specification.

\subsection{Socioeconomic gradient: HDI quintile analysis}

The linkage of DRG records to municipal HDI was successful for 87{.}3\% of
C16.* discharges. Kruskal-Wallis testing of LOS across HDI quintiles yielded
$H(4) = 28{.}7$, $p < 0{.}001$. Chi-squared testing of mortality across quintiles
yielded $\chi^2(4) = 14{.}2$, $p = 0{.}003$.

As shown in Fig.~3, the in-hospital mortality rate exhibited a monotonic gradient:
Q1 (lowest HDI municipalities): 5{.}2\%; Q2: 4{.}7\%; Q3: 4{.}1\%; Q4: 3{.}9\%;
Q5 (highest HDI): 2{.}4\%. The absolute difference between Q1 and Q5 was
2{.}8 percentage points, a difference of 117\% in relative terms. Standardized
residuals from the contingency table confirmed significant over-representation of
mortality in Q1 ($r = +2{.}7$) and under-representation in Q5 ($r = -2{.}1$).

Patients from Q1 municipalities also exhibited longer adjusted LOS
(median: 8\,days vs.\ 6\,days in Q5) and higher emergency admission rates
(50{.}3\% vs.\ 35{.}6\%), the latter serving as a proxy for late-stage presentation
and insufficient access to elective oncological care.

\subsection{Hospital typology: K-means clustering}

K-means clustering ($k = 3$, silhouette\,=\,0{.}43) identified three distinct
hospital profiles (Table~\ref{tab:kmeans}):

\begin{table}[t]
\small
\begin{tabularx}{\linewidth}{@{}L{1.8cm}C{1.2cm}C{1.2cm}C{1.2cm}C{1.2cm}@{}}
\toprule
\rowcolor{tablehead}
\textcolor{white}{\textbf{Cluster}} &
\textcolor{white}{\textbf{$n$ hosp.}} &
\textcolor{white}{\textbf{Mort.\%}} &
\textcolor{white}{\textbf{LOS}} &
\textcolor{white}{\textbf{Emerg.\%}} \\
\midrule
1 — High-complexity / Urgent &  4 & 6{.}1 & 9{.}3 & 54 \\
2 — Surgical / Specialized   &  6 & 3{.}8 & 8{.}6 & 38 \\
3 — General / Low-complexity &  5 & 2{.}9 & 5{.}8 & 41 \\
\bottomrule
\end{tabularx}
\caption{K-means cluster centroids ($k = 3$). Mean in-hospital mortality rate (\%),
mean length of stay (LOS, days), and emergency admission rate (\%).
Cluster labels are interpretive. Silhouette coefficient = 0{.}43.}
\label{tab:kmeans}
\end{table}

Cluster 1 (high-complexity/urgent) hospitals showed the highest mortality and LOS,
and the highest emergency admission rates—consistent with serving a sicker, later-stage
population. Cluster 3 (general/low-complexity) hospitals showed the lowest mean LOS
but comparable procedure counts, suggesting possible undertreatment or case-mix
differences. HDI mapping revealed that Cluster 1 hospitals disproportionately serve
municipalities in HDI quintiles Q1–Q2, reinforcing the socioeconomic–institutional
interaction identified in the preceding analyses.

% ═══════════════════════════════════════════════════
\section{Discussion}
% ═══════════════════════════════════════════════════

\subsection{Principal findings}

This study provides, to our knowledge, the first quantification of hospital-level
variation in gastric cancer outcomes across Chile's public health network using
multilevel methods and linkage to socioeconomic territorial indicators. Our
principal findings are: (i) procedural variability across hospitals is structural and
generalized, with 62\% of hospital pairs significantly different after Bonferroni
correction; (ii) procedure count is independently associated with in-hospital
mortality and LOS after controlling for GRD severity, weight, age, comorbidity, and
hospital; (iii) the treating hospital explains 18{.}4\% of the variance in LOS—a
proportion that cannot be attributed to patient characteristics; and (iv) patients
from the most socioeconomically deprived municipalities experience a mortality rate
2{.}8 percentage points higher than those from the least deprived, an excess risk
that persists after institutional adjustment.

\subsection{Comparison with international evidence}

The magnitude of hospital-level variation we observe is consistent with international
literature. Kamaraju et al.\textsuperscript{6} documented mortality differences of
up to 8 percentage points across US hospitals treating gastric cancer, after risk
adjustment. The OECD's \textit{Health at a Glance 2023}\textsuperscript{7} identifies
inter-hospital variation as a priority quality indicator across OECD members, with
procedure-based proxies accounting for 15–25\% of explained variance in LOS—a range
that brackets our ICC of 18{.}4\%.

The socioeconomic gradient in outcomes (2{.}8 pp mortality difference Q1 vs.\ Q5)
adds a dimension rarely captured in high-income country studies due to the
compressed socioeconomic range of their health systems. Our finding parallels evidence
from Colombia\textsuperscript{8} and Brazil,\textsuperscript{9} where geographic and
income-based inequalities in cancer outcomes are documented in administrative data.

\subsection{Mechanistic interpretation}

Two complementary mechanisms likely explain our findings. First, \textit{late
presentation bias}: patients from lower-HDI municipalities arrive with higher GRD
severity and more emergency presentations, which independently drive mortality and
LOS. This is captured in the statistical models through GRD severity and TIPO\_INGRESO.
Second, even after conditioning on these clinical proxies, a residual institutional
effect persists (the ICC), suggesting that differences in \textit{treatment protocols,
infrastructure, and human resources} across hospitals contribute independently to
outcomes—the classical unjustified clinical variability pathway.\textsuperscript{5}

The K-means clustering supports this dual mechanism: Cluster 1 hospitals (highest
mortality, most emergencies) serve the lowest-HDI municipalities and appear in the
BLUP forest plot as the institutions with the highest positive random intercepts.

\subsection{Class imbalance and predictive modeling}

The severely imbalanced outcome (mortality rate: 4{.}05\%) presents a practical
challenge for clinical decision support tools. At the conventional threshold of 0.5,
the baseline logistic model achieves near-zero sensitivity for the minority class. Our
comparison of three rebalancing strategies (class weighting, SMOTE oversampling) shows
that lowering the decision threshold to 0.10 substantially improves sensitivity (to
$\approx$\,60–70\%) at the cost of increased false positives—a trade-off preferable
in the clinical context of not missing high-risk admissions.

\subsection{Limitations}

This study has several limitations. The DRG database does not include cancer staging
(TNM), histology (intestinal vs.\ diffuse type), Eastern Cooperative Oncology Group
(ECOG) performance status, or specific treatment regimens (chemotherapy, surgery type).
These variables, if available, would likely explain a portion of the residual variance
(36\% unexplained in H\textsubscript{3}) and could modify the estimated ICC.

The unit of analysis is the \textit{hospitalization episode}, not the individual
patient: a patient hospitalized three times in the study period contributes three
records. While this aligns with the DRG reporting logic, it prevents individual-level
survival or trajectory analysis.

The HDI analysis operates at the municipal level, introducing the ecological fallacy
risk: we cannot infer individual-level socioeconomic status from the municipality's
HDI. Future linkage with social protection surveys (CASEN) or individual income data
would provide a more precise socioeconomic exposure.

Finally, the observational design precludes causal inference. The associations
reported reflect adjusted correlations in administrative data; confounding by
unmeasured variables (e.g., surgeon experience, hospital accreditation status)
cannot be excluded.

\subsection{Conclusions and policy implications}

Hospital-level variation in gastric cancer care in Chile is substantial, statistically
robust, and socioeconomically patterned. Three policy implications follow directly
from our evidence:

\textbf{(1) Standardize oncological protocols} in high-variability hospitals, particularly
those in Cluster 1, where the combined effect of late presentation and potentially
suboptimal institutional protocols compounds patient risk.

\textbf{(2) Invest differentially} in health establishments serving low-HDI municipalities,
addressing both preventive (reducing late diagnoses) and curative (improving institutional
capacity) dimensions of the care pathway.

\textbf{(3) Implement early referral systems} for complex cases to higher-complexity
centers, using the hospital typology identified here as a decision framework for
appropriate case routing.

These recommendations are actionable within the current FONASA governance framework
and do not require major legislative or structural changes—only evidence-based
resource allocation and protocol enforcement.

% ═══════════════════════════════════════════════════
\section*{Contributors}
% ═══════════════════════════════════════════════════
V.R., J.T.A., and S.H. contributed equally to conceptualization, data curation,
formal analysis, visualization, and writing. V.R. had primary responsibility
for statistical modeling and code development.

% ═══════════════════════════════════════════════════
\section*{Declaration of interests}
% ═══════════════════════════════════════════════════
The authors declare no competing interests. This work was carried out as a
final project for the course \textit{Análisis de Datos e Inferencia Estadística},
School of Engineering, Universidad del Desarrollo, Santiago, Chile (2026).

% ═══════════════════════════════════════════════════
\section*{Data sharing}
% ═══════════════════════════════════════════════════
The GRD database is publicly available from Chile's Ministry of Health
(\url{https://datos.gob.cl}). The analysis code, cleaned dataset, and output files
are available in the project repository at
\url{https://github.com/SebaGRT/Proyecto-Hospitalizacion}.

% ═══════════════════════════════════════════════════
\section*{Acknowledgements}
% ═══════════════════════════════════════════════════
The authors thank the Ministerio de Salud de Chile (MINSAL) and FONASA for
maintaining the publicly accessible DRG database. Computational analyses were
performed using Python 3.13 with pandas, numpy, statsmodels, scikit-learn,
matplotlib, and seaborn libraries.

% ═══════════════════════════════════════════════════
\section*{References}
% ═══════════════════════════════════════════════════
\begin{enumerate}[leftmargin=16pt,label=\arabic*,itemsep=0pt,
    parsep=0pt,topsep=2pt,font=\small]

\item Ministerio de Salud de Chile. \textit{Situación de Salud de Chile 2023:
  Indicadores Básicos}. MINSAL; 2023.

\item Global Burden of Disease Collaborative Network. \textit{Global Burden of
  Disease Study 2021}. IHME; 2022.

\item Wennberg JE. \textit{Tracking Medicine: A Researcher's Quest to Understand
  Health Care}. Oxford University Press; 2010.

\item OECD. \textit{Health at a Glance 2023: OECD Indicators}. OECD Publishing; 2023.
  \url{https://doi.org/10.1787/7a7afb35-en}

\item Wennberg JE. Unwarranted variations in healthcare delivery: implications
  for academic medical centres. \textit{BMJ}. 2002;325(7370):961–964.

\item Kamaraju S, Maganty A, Davies BJ. Hospital-level variation in the management
  and outcomes of gastric cancer. \textit{J Surg Oncol}. 2022;126(4):656–665.

\item OECD. \textit{Health at a Glance 2023}. Indicator 6.3: Avoidable hospital
  admissions. OECD Publishing; 2023.

\item Piñeros M, et al. Cancer patterns and trends in Colombia, 2000–2019.
  \textit{Lancet Reg Health Am}. 2023;17:100385.

\item de Souza DLB, et al. Geographic inequalities in cancer mortality in Brazil.
  \textit{PLoS One}. 2021;16(3):e0248580.

\item Hollander M, Wolfe DA. \textit{Nonparametric Statistical Methods}, 2nd ed.
  Wiley; 1999.

\item White H. A heteroskedasticity-consistent covariance matrix estimator and a
  direct test for heteroskedasticity. \textit{Econometrica}. 1980;48(4):817–838.

\item Bernal YA, Campaña C, Sanhueza C, et al. Multimorbidity profile among
  cancer-related hospitalization events in younger and older patients: a large-scale
  nationwide cross-sectional study. \textit{Lancet Reg Health Am}. 2026;53.

\item Sekhon M, Cartwright M, Francis JJ. Acceptability of healthcare interventions:
  an overview of reviews and development of a theoretical framework.
  \textit{BMC Health Serv Res}. 2017;17(1):88.

\item Seabold S, Perktold J. statsmodels: Econometric and statistical modeling with
  Python. \textit{Proceedings of the 9th Python in Science Conference}. 2010:92–96.

\item Pedregosa F, et al. Scikit-learn: Machine learning in Python.
  \textit{J Mach Learn Res}. 2011;12:2825–2830.

\end{enumerate}

\vspace{10pt}
\noindent\rule{\linewidth}{0.8pt}
{\fontsize{7.5}{9.5}\selectfont
Submitted: May 2026 \quad|\quad
Final version: May 2026 \quad|\quad
Universidad del Desarrollo, School of Engineering \quad|\quad
Course: Análisis de Datos e Inferencia Estadística
}

\end{document}
